\documentclass[11pt,a4paper]{article}

%% PACHETE MATEMATICE
\usepackage{amsmath,amssymb,amsfonts}
\usepackage{amsthm}
\usepackage{mathpazo}
\usepackage{eulervm}
\usepackage{enumerate}
\usepackage{color}
\usepackage{graphicx}
\usepackage[all]{xy}
\usepackage{xcolor}
\usepackage{framed}
\usepackage[unicode]{hyperref}
\setcounter{MaxMatrixCols}{20}
\usepackage{multicol}
%\usepackage[romanian]{babel}


%% ASPECT
\linespread{1.05}
\usepackage[T1]{fontenc}
\usepackage[utf8x]{inputenc}
\usepackage[margin=2cm]{geometry}
% \usepackage[color]{showkeys}
\usepackage{anyfontsize}
\usepackage{titlesec}
\titleformat*{\section}{\large\bfseries}

\usepackage{fancyhdr}
\pagestyle{fancy}
\rhead{\small \color{gray} Notițe de Adrian Manea}
\lhead{\small \color{gray} M1-ACS, 2017---2018, M. Olteanu}


\usepackage[autostyle = false, style = german]{csquotes}
\usepackage[romanian]{babel}
\newcommand\qq{\enquote}

%% COMENZILE MELE
\renewcommand*{\proofname}{Dem.:}
\newcommand\bb{\mathbb}
\newcommand\dr{\mathrm}
\newcommand\kal{\mathcal}
\newcommand\fk{\mathfrak}
\newcommand\op{\oplus}
\newcommand\ot{\otimes}
\newcommand\ds{\displaystyle}
\newcommand\id{\indent}
\newcommand\nid{\noindent}
\newcommand\seq{\subseteq}
\newcommand\sneq{\subsetneq}
\newcommand\speq{\supseteq}
\newcommand\spneq{\supsetneq}
\newcommand\rar{\rightarrow}
\newcommand\Rar{\Rightarrow}
\newcommand\xrar{\xrightarrow}
\newcommand\lar{\leftarrow}
\newcommand\Lar{\Leftarrow}
\newcommand\xlar{\xleftarrow}
\newcommand\lrar{\leftrightarrow}
\newcommand\Lrar{\Leftrightarrow}
\newcommand\ideal{\trianglelefteq}
\newcommand\ai{\text{ a.\^{i}. }}
\newcommand\sk{(\textit{Schi\c{t}\u{a}}) }
\newcommand\rhup{\rightharpoonup}
\newcommand\rhdn{\rightharpoondown}
\newcommand\lhup{\leftharpoonup}
\newcommand\lhdn{\leftharpoondown}
\newcommand\vid{\emptyset}
\newcommand\wh{\widehat}
\newcommand\ol{\overline}
\newcommand\NN{\mathbb{N}}
\newcommand\RR{\mathbb{R}}
\newcommand\ZZ{\mathbb{Z}}
\newcommand\QQ{\mathbb{Q}}
\newcommand\CC{\mathbb{C}}

%% STILURILE MELE
\newtheoremstyle{dotless}{}{}{\itshape}{}{\bfseries}{:}{ }{}

\theoremstyle{dotless}
\newtheorem{theorem}{Teorem\u{a}}[section]
\newtheorem{proposition}{Propozi\c{t}ie}[section]
\newtheorem{lemma}{Lem\u{a}}[section]
\newtheorem{corollary}{Corolar}[section]
\newtheorem{exercise}{Exerci\c{t}iu}

\newtheoremstyle{dotlessdr}{}{}{}{}{\bfseries}{:}{ }{}
\theoremstyle{dotlessdr}
\newtheorem{example}{Exemplu}[section]
\newtheorem{definition}{Defini\c{t}ie}[section]
\newtheorem{remark}{Observa\c{t}ie}[section]




\begin{document}

\begin{center}
  {\large\textbf{Seminar 7 \\ Funcții de mai multe variabile. \\ Derivate parțiale }}
\end{center}

\section{Derivate parțiale. Diferențială}
\label{sec:der-dif}

Amintim următoarele noțiuni: Fie $f : \RR^n \to \RR$ o funcție de $n$ variabile.
\begin{itemize}
\item \emph{Diferențiala (totală)} a lui $f$ se definește prin:
  \[
    Df(a)(x_1, \dots, x_n) = \sum_{i = 1}^n \frac{\partial f}{\partial x_i}(a) dx_i, \forall a \in \RR^n.
  \]
\item \emph{Diferențiala a doua} a lui $f$ se definește prin:
  \[
    D^2f(a)(x_1, \dots, x_n) = \sum_{i = 1}^n \sum_{j = 1}^n \frac{\partial^2 f}{\partial x_i \partial x_j} (a)
    \cdot dx_i dx_j, \forall a \in \RR^n.
  \]
\item \emph{Laplacianul} funcției $f$ se definește prin:
  \[
    \Delta f = \sum_{i = 1}^n \frac{\partial^2 f}{\partial x_i^2}.
  \]
\end{itemize}

\begin{definition}\label{def:armonica}
  O funcție se numește \emph{armonică} dacă laplacianul său este nul.
\end{definition}

%%%%%%%%%%%%%%%%%%%%%%%%%%%%%%%%%%%%%%%%%%%%%%%%%%%%%%%%%%%%%%%%%%%%%%

\section{Exerciții}
\label{sec:exc}

1. Verificați dacă următoarele funcții sînt armonice:
\begin{enumerate}[(a)]
\item $f : \RR^2 - \{(0, 0)\} \to \RR, f(x, y) = \ln(x^2 + y^2)$;
\item $f : \RR^3 - \{ (0, 0, 0) \} \to \RR, f(x, y, z) = \ln(x^2 + y^2 + z^2)$;
\item $f : \RR^2 - \{ (0, 0) \} \to \RR, f(x, y) = \frac{1}{\sqrt{x^2 + y^2}}$;
\item $f : \RR^3 - \{ (0, 0, 0) \} \to \RR, f(x, y, z) = \frac{1}{\sqrt{x^2 + y^2 + z^2}}$;
\item $f : \RR^2 - \{ (0, 0) \} \to \RR, f(x, y) = \frac{1}{x^2 + y^2}$;
\item $f : \RR^3 - \{ (0, 0, 0) \} \to \RR, f(x, y) = \frac{1}{x^2 + y^2 + z^2}$;
\end{enumerate}

\vspace{1cm}

2. Fie funcția:
\[
  f(x, y) = \begin{cases}
    \frac{xy^3}{x^2 + y^2}, & (x, y) \neq (0, 0) \\
    0, & (x, y) = (0, 0)
  \end{cases}
\]

\begin{enumerate}[(a)]
\item Să se arate că $f$ este de clasă $\kal{C}^1$ pe $\RR^2$;
\item Să se arate că $f$ are derivate parțiale mixte de ordinul al doilea în orice punct
  și să se calculeze $\frac{\partial^2 f}{\partial x \partial y}$ și $\frac{\partial^2 f}{\partial y \partial x}$
  în origine. Este $f$ de clasă $\kal{C}^2$ pe $\RR$?
\end{enumerate}

\vspace{1cm}

3. Să se calculeze jacobienii transformărilor în coordonate polare,
cilindrice și sferice:
\[
  \begin{cases}
    (x, y) &= (\rho \cos \varphi, \rho \sin \varphi), (\rho, \varphi) \in (0, \infty) \times (0, 2\pi) \\
    (x, y, z) &= (\rho \cos \varphi, \rho \sin \varphi, z), (\rho, \varphi, z) \in (0, \infty) \times (0, 2\pi) \times \RR \\
    (x, y, z) &= (\rho \sin \theta \cos \varphi, \rho \sin \theta \sin \varphi, \rho \cos \varphi), %
    (\rho, \theta, \varphi) \in (0, \infty) \times (0, \pi) \times (0, 2\pi)
  \end{cases}
\]

\vspace{1cm}

4. Să se calculeze diferențialele funcțiilor:
\begin{enumerate}[(a)]
\item $f : \RR^2 - \{ (x, y) \mid x = 0 \}, f(x, y) = \arctan \frac{y}{x}$;
\item $g : \RR^2 - \{ (x, y) \mid y = 0 \}, g(x, y) = -\arctan \frac{x}{y}$.
\end{enumerate}

\vspace{1cm}

5. Fie $f \in \kal{C}^2(\RR^2)$ și $g : \RR^2 \to \RR$, dată de $g(x, y) = f(x^2 + y^2, x^2 - y^2)$.
Să se calculeze derivatele parțiale de ordinul întîi și al doilea ale funcției $g$, precum și
diferențiala acestei funcții.

\vspace{1cm}

6. Fie $r = \sqrt{x^2 + y^2 + r^2}$ și fie $f \in \kal{C}^1(\RR)$. Să se calculeze laplacianul
funcțiilor $g(x, y, z) = f \big( \frac{1}{r} \big)$ și $h(x, y, z) = f(r)$.

\vspace{1cm}

7. Dacă $f \in \kal{C}^2(\RR^3)$ și $u(x, y) = f(x^2 + y^2, x^2 - y^2, 2xy)$, să se calculeze
derivatele parțiale de ordinul al doilea ale funcției $u$.

\vspace{1cm}

8. Fie $a \in \RR$ și $g, h$ două funcții de clasă $\kal{C}^2$ pe $\RR$. Să se arate că
$f(x, y) = g(x - ay) + h(x + ay)$ verifică ecuația coardei vibrante:
\[
  \frac{\partial^2 f}{\partial x^2} - a^2 \frac{\partial^2 f}{\partial y^2} = 0.
\]

\vspace{1cm}

9. Să se afle $f \in \kal{C}^2(\RR)$, știind că funcția $v(x, y) = f\big(\frac{y}{x}\big)$
este armonică.

\vspace{1cm}

10. Să se arate că laplacianul în coordonate cilindrice este:
\[
  \Delta f = \frac{1}{\rho} \frac{\partial}{\partial \rho} %
  \Big(\rho \frac{\partial f}{\partial \rho} \Big) + %
  \frac{1}{\rho^2} \frac{\partial^2 f}{\partial \varphi^2} + %
  \frac{\partial^2 f}{\partial z^2}.
\]

Indicație: porniți cu definiția laplacianului în coordonate carteziene, apoi
folosiți formula pentru derivarea funcțiilor compuse, ținînd cont de formulele
de trecere în coordonate cilindrice. De exemplu, pentru a calcula
$\frac{\partial f}{\partial \rho}$, folosiți
$\frac{\partial f}{\partial x} \cdot \frac{\partial x}{\partial \rho}$ etc.

\vspace{1cm}

11. Să se arate că funcția $z(x, y) = xy + xe^{\frac{y}{x}}$ verifică ecuația:
\[
  x \frac{\partial z}{\partial x} + \frac{\partial z}{\partial y} = xy + z.
\]
\end{document}
